\documentclass[12pt,a4paper]{ctexart}

% ── 编码与字体 ──
\usepackage{amsmath,amssymb}

% ── 页面设置 ──
\usepackage[margin=2.5cm]{geometry}
\usepackage{setspace}
\onehalfspacing

% ── 插图与表格 ──
\usepackage{graphicx}
\usepackage{float}
\usepackage{caption}
\usepackage{booktabs}
\usepackage{array}

% ── 颜色与链接 ──
\usepackage[colorlinks=true,linkcolor=blue,urlcolor=blue]{hyperref}

% ── 代码高亮 ──
\usepackage{listings}
\usepackage{xcolor}

\lstset{
    language=C++,
    basicstyle=\small\ttfamily,
    backgroundcolor=\color{gray!5},
    frame=single,
    breaklines=true,
    numbers=left,
    numberstyle=\tiny,
    keywordstyle=\color{blue},
    commentstyle=\color{green!50!black},
    stringstyle=\color{red!50!orange},
    tabsize=4,
    captionpos=b
}

\lstdefinestyle{csharp}{
    language=[Sharp]C,
    basicstyle=\small\ttfamily,
    backgroundcolor=\color{gray!5},
    frame=single,
    breaklines=true,
    numbers=left,
    numberstyle=\tiny,
    keywordstyle=\color{blue},
    commentstyle=\color{green!50!black},
    stringstyle=\color{red!50!orange},
    tabsize=4,
    captionpos=b
}

\title{\vspace{-2cm}Desktop Pet --- 桌面宠物技术文档}
\author{符玄桌宠项目}
\date{\today}

\begin{document}

\maketitle
\thispagestyle{empty}

\begin{abstract}
本文档为桌面宠物项目的完整技术文档。项目基于 Unity 2022.3 LTS 引擎，
使用 Live2D Cubism SDK 5-r.4 渲染《崩坏：星穹铁道》角色「符玄」，
通过 Win32 透明窗口叠加技术实现桌面宠物效果。
文档涵盖系统架构、各模块设计、动画系统、状态机及交互机制。
\end{abstract}

\tableofcontents
\newpage

% ═══════════════════════════════════════════
\section{项目概述}
% ═══════════════════════════════════════════

\subsection{项目背景}

原桌面宠物程序基于 Windows GDI+ 分层窗口技术（\texttt{text1.c}），使用静态 PNG 图片硬切换模拟动画。
本项目将其迁移至 Unity 引擎，利用 Live2D Cubism SDK 实现参数化变形动画，
获得更流畅的视觉效果和更丰富的交互能力。

\subsection{技术栈}

\begin{itemize}
    \item \textbf{引擎}：Unity 2022.3.62t7 (LTS)
    \item \textbf{角色渲染}：Live2D Cubism SDK 5-r.4（参数化变形 + 实时物理模拟）
    \item \textbf{窗口管理}：Win32 API（user32.dll）--- 色键透明、置顶、点击穿透
    \item \textbf{交互方式}：左键拖拽/抛掷、右键上下文菜单
    \item \textbf{语言}：C\# 脚本
    \item \textbf{平台}：Windows 10/11（Win64）
\end{itemize}

\subsection{项目状态}

功能已完成：
\begin{itemize}
    \item Live2D 模型加载与渲染（自动 AssetDatabase / Resources 双路径）
    \item 桌面透明窗口（色键抠除 \#00FF00 + WS\_EX\_LAYERED）
    \item 鼠标穿透动态管理（WS\_EX\_TRANSPARENT）
    \item 像素物理引擎（重力、边界碰撞、抛掷）
    \item 地面任务状态机（5 种行为加权随机切换）
    \item 走路动画（侧身转体 + 腿臂交叉摆动 + 上下颠簸 + 体态淡入淡出）
    \item 11 种空闲动作循环（歪头、微笑、挑眉、星辉、伸懒腰、爱心、数钱、委屈、法阵、害羞、困惑）
    \item 右键上下文菜单（权重调整 + 强制动作播放）
    \item 强制动作锁定机制（播放期间不被走路覆盖，播完自动恢复）
    \item 分区点击反馈（头部/身体/腿部 3 区不同反应）
    \item 拖拽挣扎动画（双臂交替划水 + 双腿交替 + 身体扭动 + 表情）
    \item 头发/裙子/法盘直接物理驱动（拖拽速度实时输入 CubismPhysics）
    \item 时间/天气响应系统（昼夜犯困眼皮 + 天气表情基调 + 待机气泡）
    \item AI 对话 + 自动闲聊 + 底部输入栏
    \item 调试窗口实时调参
    \item 3D 模型渲染管线预留（HybridRenderer 双渲染器交叉淡入淡出）
\end{itemize}

% ═══════════════════════════════════════════
\section{系统架构}
% ═══════════════════════════════════════════

\subsection{整体架构}

\begin{lstlisting}[frame=single,numbers=none,style=csharp]
+-----------------------------------------------------------------+
|            Unity 窗口 (透明叠加层)                                |
|  +-----------------------------------------------------------+  |
|  |  Camera (Clear=SolidColor, BG=#00FF00)                     |  |
|  |                                                              |  |
|  |  DesktopPet (主控)                                          |  |
|  |    +-- 像素物理 (重力/碰撞/落地)                             |  |
|  |    +-- 地面任务状态机 (加权随机选择)                          |  |
|  |                                                              |  |
|  |  DragHandler (交互)                                         |  |
|  |    +-- 左键拖拽窗口移动 + 抛掷                               |  |
|  |    +-- 右键菜单开关                                          |  |
|  |    +-- 点击穿透管理 (每帧)                                   |  |
|  |                                                              |  |
|  |  Live2DRenderer (渲染)                                      |  |
|  |    +-- CubismModel 参数驱动                                  |  |
|  |    +-- 眨眼/呼吸/噪声微动                                    |  |
|  |    +-- 空闲动作 1-10 循环                                    |  |
|  |    +-- 走路动画 (转体/抬腿/颠簸)                            |  |
|  |    +-- 体态淡入淡出 (0.3s/0.4s)                             |  |
|  |    +-- 强制动作锁定 (actionLocked)                           |  |
|  |                                                              |  |
|  |  ContextMenu (UI)                                           |  |
|  |    +-- 权重调整 (5 种任务权重)                              |  |
|  |    +-- 强制动作触发按钮 (10 个动作)                          |  |
|  |    +-- OnGUI 直接绘制 (无 Canvas)                            |  |
|  |                                                              |  |
|  |  WindowOverlay (Win32)                                      |  |
|  |    +-- SetWindowLong (透明/穿透)                             |  |
|  |    +-- SetLayeredWindowAttributes (色键)                     |  |
|  |    +-- SetWindowPos (全屏置顶)                               |  |
|  +-----------------------------------------------------------+  |
+-----------------------------------------------------------------+
\end{lstlisting}

\subsection{组件依赖关系}

关键交互流程：

\begin{description}
    \item[渲染循环] DesktopPet.Update $\to$ Live2DRenderer.LateUpdate（参数设置）
    \item[物理更新] Update $\to$ StepPet（重力/速度/位置）$\to$ UpdateGroundTask（状态机）
    \item[强制动作] ContextMenu.DrawActionButton $\to$ DesktopPet.ForceStop $\to$ \\[2pt]
        \hspace*{2em}Live2DRenderer.ForceIdleAction $\to$（锁定）$\to$ ResetIdleAction（解锁+回调）
    \item[菜单交互] DragHandler.Update（Block 0: 菜单穿透管理）$\to$ ContextMenu.IsOpen
    \item[走路淡入] LateUpdate 检测 isWalking 切换 $\to$ 启动 walkFadeInRemaining（0.3s 体态渐入）
\end{description}

% ═══════════════════════════════════════════
\section{窗口管理 (WindowOverlay)}
% ═══════════════════════════════════════════

\subsection{色键透明原理}

\begin{enumerate}
    \item Camera 设置 ClearFlags = SolidColor，Background = (0, 1, 0) 即 \#00FF00 亮绿色
    \item 通过 Win32 API \texttt{SetWindowLong} 设置 \texttt{WS\_EX\_LAYERED} 扩展样式
    \item \texttt{SetLayeredWindowAttributes} 将颜色键 0x0000FF00 设为透明
    \item 构建版自动重试 5 次查找窗口句柄（延迟帧防竞态）
\end{enumerate}

\subsection{点击穿透 (Click-Through)}

\begin{itemize}
    \item \texttt{WS\_EX\_TRANSPARENT} 样式让鼠标事件透过透明区域传到桌面
    \item 每帧根据鼠标位置动态切换：鼠标在宠物区域内 $\to$ 移除穿透（接收点击）；
        鼠标在区域外 $\to$ 添加穿透（点击不受影响）
    \item 菜单打开时：菜单区域内可交互，区域外穿透（优先于普通穿透逻辑）
\end{itemize}

\subsection{窗口样式}

\begin{itemize}
    \item 全屏无边框窗口（WS\_POPUP）
    \item 置顶（HWND\_TOPMOST）
    \item 不在任务栏显示（WS\_EX\_TOOLWINDOW）
    \item 工具窗口风格
\end{itemize}

% ═══════════════════════════════════════════
\section{物理引擎 (DesktopPet)}
% ═══════════════════════════════════════════

\subsection{坐标系}

采用像素坐标系：左上角为原点 (0, 0)，X 轴向右为正，Y 轴向下为正。
物理状态全部用整数像素值表示，避免浮点精度问题。

\subsection{物理参数}

\begin{table}[H]
\centering
\caption{核心物理常量}
\label{tab:physics}
\begin{tabular}{lc}
\toprule
\textbf{参数} & \textbf{值} \\
\midrule
重力加速度 & 1 px/frame² \\
最大下落速度 & 8 px/frame \\
最大水平速度 & 12 px/frame \\
地面 Y 偏移 & -300 px（相对于屏幕底部） \\
速度系数 & 0.5 \\
\bottomrule
\end{tabular}
\end{table}

\subsection{地面任务状态机}

宠物共有 5 种地面行为，通过加权随机选择：

\begin{table}[H]
\centering
\caption{GroundTask 状态机}
\label{tab:tasks}
\begin{tabular}{lcl}
\toprule
\textbf{任务} & \textbf{权重(默认)} & \textbf{行为} \\
\midrule
MoveLeftEdge  & 1   & 向左走到屏幕左边缘 \\
MoveRightEdge & 1   & 向右走到屏幕右边缘 \\
MoveLeftTime  & 1   & 向左走随机时长 (2-4s) \\
MoveRightTime & 1   & 向右走随机时长 (2-4s) \\
StopTime      & 6   & 停止随机时长 (6-15s) \\
\bottomrule
\end{tabular}
\end{table}

状态机特点：
\begin{itemize}
    \item 到达边缘后自动转向：左边缘 $\to$ 右走/停止；右边缘 $\to$ 左走/停止
    \item 定时任务到期 $\to$ 重新加权随机选择
    \item 停止任务期间播放空闲动作循环
\end{itemize}

\subsection{暂停/恢复机制}

\begin{itemize}
    \item \texttt{isPaused} 布尔值：暂停时物理 Update 直接 return
    \item \texttt{ForceStop()}：将边缘任务转为定时任务，清零 petVx，启动停止任务
    \item \texttt{Resume()}：恢复时如果当前无任务或停止中，启动新任务
    \item \texttt{Pause(float)}：暂停指定秒数后自动恢复
\end{itemize}

% ═══════════════════════════════════════════
\section{拖拽挣扎动画 (Live2DRenderer)}
% ═══════════════════════════════════════════

\subsection{设计目标}

拖拽时角色应表现出\"被提着走\"的挣扎感，同时裙子和头发有惯性拖尾效果。
物理系统（CubismPhysicsController）在拖拽中继续运作，但参数值会被挣扎动画覆盖。

\subsection{核心机制}

拖拽挣扎动画通过帧间速度驱动：每帧计算 $\Delta x = petX_{t} - petX_{t-1}$，
经平滑滤波后驱动各参数：

\begin{itemize}
    \item \textbf{双臂交替划水}：右臂正=向前，左臂负=向前，同相位实现交替
        \begin{itemize}
            \item 大幅摆动：Param94（上臂旋转，范围 [-30,60]）
            \item 小幅：Param31/32/33（右臂），Param34/36/37（左臂）
            \item 手部透视图层跟随幅度（自动浮到衣服前）
        \end{itemize}
    \item \textbf{双腿交替}：Param126/129 腿前后摆动，Param127/131 腿弯曲
    \item \textbf{身体扭动}：ParamBodyAngleY 跟随鼠标方向平滑转身
    \item \textbf{表情}：瞪眼 + 张嘴（\"慌张\"表情）
\end{itemize}

\subsection{直接物理驱动}

帧间速度经平滑后直接写入 Cubism 参数，为 CubismPhysics 提供输入信号：

\begin{itemize}
    \item \textbf{裙子}：Param82（前摆动）、Param87（后摆动）、Param84（侧摆动）
    \item \textbf{法盘}：Param49/51/57/60（法盘旋转/摆动）
    \item \textbf{头发}：20 个输出参数（Param5/7/9/11/14/17/19/21/23/35/41/43/45/55/62/91/74/89/169 等）
    \item \textbf{头部}：ParamAngleX/Z 跟随速度方向
\end{itemize}

所有直接驱动参数绕过物理系统的 Delay=0.8s 延迟，实现即时响应。

% ═══════════════════════════════════════════
\section{时间/天气响应系统 (TimeWeatherController)}
% ═══════════════════════════════════════════

\subsection{概述}

\texttt{TimeWeatherController} 负责感知系统时间和当地天气，向 \texttt{Live2DRenderer} 提供状态信息，使宠物的表情、动作权重和待机气泡内容随昼夜/天气变化。

\subsection{时间系统}

\begin{itemize}
    \item 使用 \texttt{System.DateTime.Now} 获取当地小时（0-23）
    \item \texttt{isNight}：18:00 - 6:00（夜间模式）
    \item \texttt{isSleepyTime}：22:00 - 5:00（犯困时段）
    \item \texttt{dayPhase}：6:00 = 0.0，12:00 = 0.5，18:00 = 1.0（用于渐变）
\end{itemize}

\subsection{天气系统}

\begin{itemize}
    \item 通过 wttr.in API 获取当地天气（城市可配，留空=自动 IP 定位）
    \item 轮询间隔默认 300 秒（可配）
    \item 解析天气类型：Clear / Cloudy / Overcast / Rain / Drizzle / Thunder / Snow / Fog
    \item 同时解析温度（用于气泡文本）
    \item 首次启动立即获取，失败时保留 Unknown 状态
\end{itemize}

\subsection{表情联动}

在 \texttt{UpdateIdleAnimation()} 中叠加昼夜/天气基调表情：

\begin{table}[H]
\centering
\caption{昼夜/天气表情联动}
\label{tab:weather-emotion}
\begin{tabular}{ll}
\toprule
\textbf{条件} & \textbf{表情} \\
\midrule
犯困时段（22-5点） & 眼皮微垂（EyeOpen lerp 0.7） \\
夜晚（18-22点）    & 轻度眼皮下垂 \\
阴雨/雷雨/阴天    & 眉毛微抬 + 嘴巴微嘟（委屈） \\
晴/多云          & 自然微笑 \\
雪              & 微微张嘴（好奇）+ 眼睛睁大 \\
\bottomrule
\end{tabular}
\end{table}

\subsection{动作权重调节}

\texttt{PickWeightedIdleAction()} 动态调整权重：

\begin{itemize}
    \item 夜间：星辉/爱心眼减少，委屈增加
    \item 犯困时段：委屈+2，歪头+1
    \item 雨天：委屈+2，爱心眼/数钱减少
    \item 晴天：爱心眼+1，星辉+1
    \item 雪天：挑眉+1，委屈+1
\end{itemize}

\subsection{待机气泡联动}

\texttt{PickIdleBubbleMessage()} 50\% 概率触发时间/天气特化气泡：

\begin{itemize}
    \item 犯困时段：\"好晚了…该睡了~\"\"眼睛快睁不开了…\"
    \item 夜晚：\"晚上好呀~\"\"月色真美…\"
    \item 早晨（5-8点）：\"早安~\"\"今天也要加油哦！\"
    \item 雨天：\"下雨了呢~\"\"听雨声发呆~\"
    \item 雷雨：\"打雷了好可怕！\"\"轰隆隆…好吓人…\"
    \item 雪：\"下雪了！好漂亮~\"\"想堆雪人…\"
    \item 晴天：\"今天天气真好~\"\"太阳暖洋洋的~\"
\end{itemize}

% ═══════════════════════════════════════════
\section{Live2D 渲染器 (Live2DRenderer)}
% ═══════════════════════════════════════════

\subsection{概述}

\texttt{Live2DRenderer} 是核心渲染组件，通过 Cubism SDK 驱动模型参数实现所有动画。
继承自 \texttt{IPetRenderer} 接口，被 \texttt{HybridRenderer} 包装管理。

\subsection{生命周期}

\begin{enumerate}
    \item Start: 自动加载模型（优先 AssetDatabase，降级 Resources.Load）
    \item Update: 计算走路相位 + 垂直颠簸偏移 + 模型位置同步
    \item LateUpdate: 设置所有参数（眨眼、呼吸、噪声、空闲动作/走路）
\end{enumerate}

\subsection{参数驱动系统}

使用 \texttt{SetParameter(name, value)} 统一设置 Cubism 参数，
涵盖约 60+ 个参数 ID：
\begin{itemize}
    \item \textbf{身体}：ParamBodyAngleX/Y/Z、ParamBreath
    \item \textbf{头部}：ParamAngleX/Y/Z、ParamEyeBallX/Y
    \item \textbf{眼睛}：ParamEyeLOpen/ROpen、ParamEyeLSmile/RSmile
    \item \textbf{嘴部}：ParamMouthForm
    \item \textbf{眉毛}：ParamBrowRY/LY
    \item \textbf{特效}：Param121（黑幕切换）、Param132（眼镜发光）、
        Param431/911/741（星辉环显隐）等
\end{itemize}

\subsection{空闲动作系统}

共有 10 种空闲动作循环播放（1 $\to$ 2 $\to$ $\cdots$ $\to$ 10 $\to$ 1）：

\begin{table}[H]
\centering
\caption{空闲动作列表}
\label{tab:idle}
\begin{tabular}{clc}
\toprule
\textbf{ID} & \textbf{名称} & \textbf{时长(s)} \\
\midrule
1 & 歪头 (Tilt)    & 2.0 \\
2 & 微笑 (Smile)   & 2.0 \\
3 & 挑眉 (Brow)    & 2.0 \\
4 & 星辉 (StarSpin) & 6.0 \\
5 & 伸懒腰 (Stretch) & 4.5 \\
6 & 爱心 (HeartEyes) & 3.0 \\
7 & 数钱 (Money)   & 3.5 \\
8 & 委屈 (Cry)     & 3.5 \\
9 & 法阵 (MagicCircle) & 5.5 \\
10 & 害羞 (Blush)  & 3.5 \\
\bottomrule
\end{tabular}
\end{table}

每个动作使用缓入缓出曲线 $f(t)=\sin(t\pi)$，其中 $t\in[0,1]$。
动作结束时调用 \texttt{ResetIdleAction()} 清理参数。

\textbf{选择策略}：每次播放完，通过加权随机选取下一个。各动作权重可调：

\begin{table}[H]
\centering
\caption{空闲动作权重}
\label{tab:idle-weights}
\begin{tabular}{lcc}
\toprule
\textbf{动作} & \textbf{权重} & \textbf{相对概率} \\
\midrule
歪头   & 5 & 17.9\% \\
微笑   & 5 & 17.9\% \\
挑眉   & 3 & 10.7\% \\
星辉   & 2 & 7.1\% \\
伸懒腰 & 2 & 7.1\% \\
爱心   & 3 & 10.7\% \\
数钱   & 2 & 7.1\% \\
委屈   & 3 & 10.7\% \\
法阵   & 1 & 3.6\% \\
害羞   & 3 & 10.7\% \\
\bottomrule
\end{tabular}
\end{table}

\textbf{暂停时}：\texttt{isPaused = true} 时不播新动作（菜单打开时冻结动作循环）。

\subsection{走路动画}

采用横版侧面走路风格，所有参数在 LateUpdate 中设置，利用 Cubism SDK 的 ICubismUpdatable 执行顺序机制确保物理与动画正确协作：

\begin{itemize}
    \item \textbf{转体}：ParamBodyAngleY = 18°（身体转向侧面，方向自动匹配翻转）
    \item \textbf{前倾}：ParamBodyAngleX = 5°
    \item \textbf{低头}：ParamAngleY = 8°
    \item \textbf{步频}：5 Hz（\texttt{WALK\_SWAY\_FREQ}）
    \item \textbf{抬腿}：Param165（左）/ Param164（右），幅度 4
    \item \textbf{摆臂}：交叉对位（左腿前 $\to$ 右手前），分大小范围参数：
        \begin{itemize}
            \item \textbf{大范围}：Param94（右手上臂旋转，范围 [-30,60]），$\text{legPhase}\times 2$
            \item \textbf{小范围}：Param31/32/33（右臂，范围 [-1,1]），$\text{legPhase}\times 0.4$
            \item \textbf{左臂}：Param34/36/37（左臂，范围 [-1,1]），$\text{rightPhase}\times 0.4$
        \end{itemize}
    \item \textbf{颠簸}：$1 - |\sin(\text{phase})|\times 4\,\text{px}$（迈步时最低）
    \item \textbf{呼吸加深}：ParamBreath = 3 + 相位波
\end{itemize}

\subsubsection{执行顺序与 Cubism Physics 协作}

Cubism SDK 通过 \texttt{CubismUpdateController} 统一调度所有 \texttt{ICubismUpdatable} 组件，
各组件按 \texttt{ExecutionOrder} 顺序执行：

\begin{table}[H]
\centering
\caption{Cubism 更新执行顺序}
\label{tab:exec-order}
\begin{tabular}{p{0.35\textwidth}cc}
\toprule
\textbf{组件} & \textbf{ExecutionOrder} & \textbf{功能} \\
\midrule
\texttt{Live2DRenderer.Update} & Unity Update & 设体态参数 \\
\texttt{CubismPhysicsController} & 800 & 物理模拟驱动衣服飘动 \\
\texttt{Live2DRenderer.LateUpdate} & 801 & 设腿/臂摆动 + ForceUpdateNow \\
\texttt{CubismRenderController} & 10000 & 面向相机 \\
\bottomrule
\end{tabular}
\end{table}

\textbf{关键设计}：体态参数（身体转体/前倾/低头）\\
在 \texttt{Update()} 中提前设置，使 \texttt{CubismPhysicsController(800)} \\
在 LateUpdate 阶段能读取到正确的走路体态来计算衣服飘动。\\
手臂参数在 \texttt{LateUpdate(801)} 中设置（晚于物理），\\
再通过 \texttt{ForceUpdateNow()} 强制网格重算，\\
确保手臂摆动值覆盖物理输出，同时衣服飘动由物理驱动保持自然。

\textbf{常见问题}：
\begin{enumerate}
    \item 不加 \texttt{[DefaultExecutionOrder(801)]} 时，\texttt{LateUpdate} 与 Cubism 更新顺序不确定，
        可能导致物理覆盖手臂参数（手臂值被物理置零）
    \item 不加 \texttt{ForceUpdateNow()} 时，网格使用 Update 阶段的旧参数计算，
        手臂参数虽已设但网格不变（视觉上不生效）
    \item 体态不在 Update 中提前设时，物理使用空闲体态驱动衣服，
        走路时衣服呈空闲态悬挂（不随身体飘动）
\end{enumerate}

\subsection{体态过渡}

\begin{table}[H]
\centering
\caption{过渡参数}
\label{tab:blend}
\begin{tabular}{lcc}
\toprule
\textbf{过渡} & \textbf{时长} & \textbf{曲线} \\
\midrule
走路 $\to$ 空闲（消退） & 0.4s & 二次缓出 \\
空闲 $\to$ 走路（淡入） & 0.3s & 二次缓入 \\
\bottomrule
\end{tabular}
\end{table}

走路 $\to$ 空闲时，走路的体态参数（转体/前倾/低头）经 0.4s 渐消，
同时空闲参数已经开始播放，实现平滑过渡。
空闲 $\to$ 走路时，腿臂立刻摆动（走路感瞬间有），体态经 0.3s 渐入（不走硬边）。

\subsection{强制动作锁定机制}

右键菜单触发强制动作时的保护机制：

\begin{enumerate}
    \item \texttt{ForceIdleAction(id)} 设置 \texttt{\_actionLocked = true}
    \item \texttt{isWalking} 检测条件中加入 \texttt{!\ \_actionLocked}，锁定期间不切走路
    \item 动作播放完毕 $\to$ \texttt{ResetIdleAction()} 解锁 + 触发 \texttt{OnForcedActionFinished} 事件
    \item 回调恢复宠物运动（取消暂停 + Resume）
\end{enumerate}

% ═══════════════════════════════════════════
\section{交互系统 (DragHandler)}
% ═══════════════════════════════════════════

\subsection{左键拖拽}

\begin{enumerate}
    \item 鼠标在宠物区域内按下左键 $\to$ 记录拖拽起点
    \item 移动超过 5px 阈值进入拖拽模式
    \item 拖拽中：SetWindowPos 跟随鼠标移动，关闭穿透
    \item 释放时：计算末速度（缓冲滑动平均），施加给宠物（保证 vy > 2）
    \item 进入自由落体状态
\end{enumerate}

\subsection{分区点击反馈}

点击时根据点击位置（归一化高度 0-1）分 3 区反馈：

\begin{itemize}
    \item \textbf{头部区域}（hitNormY $<$ 0.35）：戳脸 — 睁大眼 + 张嘴 + 挑眉，如"呜哇！戳到脸了"
    \item \textbf{身体区域}（hitNormY $<$ 0.65）：摸身体 — 害羞低头 + 微笑 + 眼神躲闪
    \item \textbf{腿部区域}（hitNormY $\geq$ 0.65）：踢腿 — 腿向前伸 + 坏笑 + 眯眼
\end{itemize}

反应持续期间设置 \texttt{\_poseLocked = true}，参数保存在 \texttt{\_clickSavedParams} 字典中，\texttt{LateUpdate} 重新写入以对抗 \texttt{CubismPhysicsController} 的覆盖。

\subsection{右键菜单}

\begin{enumerate}
    \item 鼠标在宠物区域内按下右键 $\to$ 打开 ContextMenu
    \item 菜单打开时：区域内可交互，区域外点击穿透
    \item 右键再次点击可关闭菜单
    \item 点击动作按钮：ForceStop $\to$ ForceIdleAction $\to$ 菜单关 UI $\to$ 动作播完自动恢复
\end{enumerate}

\subsection{点击穿透管理}

每帧执行（DragHandler.Update）：
\begin{enumerate}
    \item 菜单打开时：优先菜单穿透规则
    \item 菜单关闭时：鼠标在宠物矩形内 $\to$ 关穿透；鼠标在外 $\to$ 开穿透
    \item 确保猫爪（WS\_EX\_TRANSPARENT）不吞掉输入事件
\end{enumerate}

% ═══════════════════════════════════════════
\section{上下文菜单 (ContextMenu)}
% ═══════════════════════════════════════════

\subsection{UI 实现}

使用 Unity 的 \texttt{OnGUI} 系统直接绘制（无需 Canvas / UI Prefab），
适用于透明窗口架构。

样式特点：
\begin{itemize}
    \item 深色 VS Code 风格背景（0.15, 0.15, 0.17, 0.95）
    \item 粉色标题（0.9, 0.6, 0.8）+ 粉色强调色（0.8, 0.4, 0.6）
    \item 半透明圆角矩形背景
    \item 可滚动内容区域
\end{itemize}

\subsection{功能区域}

\begin{enumerate}
    \item \textbf{权重设置}：5 种地面任务的权重滑块（0-10），应用按钮即时生效
    \item \textbf{强制动作}：10 个动作按钮，点击后强制播放对应动作
    \item \textbf{关闭按钮}：恢复宠物运动 + 清理回调
\end{enumerate}

\subsection{动作执行流程}

点击动作按钮时的完整流程：

\begin{enumerate}
    \item \texttt{ForceStop()} $\to$ 停止走路，清零 petVx，切到 StopTime
    \item \texttt{ForceIdleAction(id)} $\to$ 设置 \texttt{\_actionLocked = true} + 播放动作
    \item \texttt{\_isOpen = false} $\to$ 菜单 UI 关闭（OnGUI 不再绘制）
    \item 注册 \texttt{OnForcedActionFinished} 回调
    \item 动作自然播完（easing 曲线结束）$\to$ \texttt{ResetIdleAction()} 触发回调
    \item 回调：取消暂停 + Resume $\to$ 宠物自动开始走路/停止
    \item 走路体态经 0.3s 淡入，过渡平滑
\end{enumerate}

% ═══════════════════════════════════════════
\section{混合渲染 (HybridRenderer)}
% ═══════════════════════════════════════════

\subsection{设计目标}

为未来引入 3D 模型渲染预留的混合渲染层。
当前仅 Live2D 渲染器活跃，3D 渲染器已挂载但 disabled。

\subsection{切换规则}

\begin{itemize}
    \item 空闲/点击/停止 $\to$ Live2D（精细表情 + 物理飘动）
    \item 走路/飞行/空中 $\to$ 3D（骨骼动画，future work）
    \item 过渡方式：Alpha 交叉淡入淡出（0.3s）
\end{itemize}

\subsection{接口设计}

\texttt{IPetRenderer} 接口定义了通用渲染器契约：

\begin{lstlisting}[style=csharp,caption={IPetRenderer 接口},label={lst:ipet}]
public interface IPetRenderer
{
    void ShowDragPose();
    void ShowClickPose();
    void ShowLandPose();
    void ShowWalkPose();
    void ShowStopPose(float lockSeconds);
    void OnPetUpdate(int petX, int petY, int petWidth, int petHeight,
                     int petVx, int petVy, bool onGround,
                     bool isDragging, bool isPaused);
}
\end{lstlisting}

% ═══════════════════════════════════════════
\section{动画详述}
% ═══════════════════════════════════════════

\subsection{空闲动作详解}

\subsubsection{动作 1 --- 歪头}
\begin{itemize}
    \item 使用 $\sin(t\pi/2)$ 缓入缓出曲线
    \item ParamAngleZ = 8°（歪头）、2s 完成
\end{itemize}

\subsubsection{动作 2 --- 微笑}
\begin{itemize}
    \item ParamEyeLSmile / RSmile = 0.6（眯眼笑）
    \item ParamMouthForm = 0.4（张嘴微笑）
    \item 2s 完成
\end{itemize}

\subsubsection{动作 3 --- 挑眉}
\begin{itemize}
    \item ParamBrowRY / LY = 6°（眉毛抬起）
    \item 2s 完成
\end{itemize}

\subsubsection{动作 4 --- 星辉环绕}
复杂多参数协同动画：
\begin{itemize}
    \item 三环（内/中/外紫环）不同的旋转速度，缓入缓出
    \item 黑幕背景（Param121 + Param137）+ 眼镜发光（Param132）
    \item 星星闪烁 + 外围星脉动
    \item 6s 完成
\end{itemize}

\subsubsection{动作 5 --- 伸懒腰}
\begin{itemize}
    \item 右手切换 + 伸出（Param92 + Param118 + 手指张合）
    \item 身体后仰（ParamBodyAngleX = $-5^\circ$）
    \item 眯眼 + 张嘴（打哈欠表情）+ 深吸一口气
    \item 4.5s 梯形曲线（快起 $\to$ 保持 $\to$ 快落）
\end{itemize}

\subsubsection{动作 6 --- 爱心眼}
\begin{itemize}
    \item 爱心眼闪烁（Param109 = 爱心眼开关）
    \item 歪头 12° + 微笑 + 身体微倾
    \item 3s 完成
\end{itemize}

\subsubsection{动作 7 --- 数钱}
\begin{itemize}
    \item 钱表情（Param122）+ 双眼放光（Param121/137/132）
    \item 美滋滋微笑 + 歪头 + 身体摇晃
    \item 3.5s 完成
\end{itemize}

\subsubsection{动作 8 --- 委屈}
\begin{itemize}
    \item 泪眼（Param130）+ 低头（ParamAngleY = 6°）
    \item 嘴巴微颤 + 八字眉（眉头抬起）
    \item 身体微微抽动（抽泣感）
    \item 3.5s 完成
\end{itemize}

\subsubsection{动作 9 --- 法阵显现}
最复杂的多阶段动画（4 个 Phase）：
\begin{itemize}
    \item Phase 1 (0-35\%)：黑幕开启 + 七星盘浮现
    \item Phase 2 (20-50\%)：白圈展开 + 呼吸位移
    \item Phase 3 (35-65\%)：眼镜发光 + 镜头微动 + 人物缩小
    \item Phase 4 (75-100\%)：各项参数二次缓出消退
    \item 5.5s 完成
\end{itemize}

\subsubsection{动作 10 --- 害羞}
\begin{itemize}
    \item 黑脸（Param101）+ 低头 + 眼神躲闪
    \item 害羞微笑 + 眯眼 + 身体扭捏
    \item 3.5s 完成
\end{itemize}

\subsubsection{动作 11 --- 困惑（AI 触发专用）}
\begin{itemize}
    \item 歪头 + 皱眉 + 眯眼
    \item 权重为 0（不自发触发，仅外部强制调用）
    \item 3s 完成
\end{itemize}

\subsection{眨眼系统}

\begin{itemize}
    \item 间隔随机 2-5s 一次
    \item 使用 $|\sin(20t)|$ 快速闭眼睁眼（0.15s）
    \item 通过 ParamEyeLOpen / ROpen 控制
\end{itemize}

\subsection{呼吸与微动}

使用 Perlin 噪声实现自然微动：
\begin{itemize}
    \item 呼吸：ParamBreath = PerlinNoise(phase, 0) $\times$ 0.4
    \item 身体晃动：3 轴 Perlin 噪声 $\times$ 幅度常量
    \item 头部微动：2 轴 Perlin 噪声 $\times$ 0.6/0.4
    \item 眼球浮动：2 轴 Perlin 噪声 $\times$ 3/2
\end{itemize}

% ═══════════════════════════════════════════
\section{目录结构}
% ═══════════════════════════════════════════

\begin{lstlisting}[frame=single,numbers=none,style=csharp]
Desktop_per_pro/
+-- README.md                              # 项目说明
+-- project_brief/
|   +-- report.pdf                         # 本文档 (PDF)
|   +-- report.tex                         # 本文档 (LaTeX)
+-- code/desktop unity/
    +-- desktop unity.sln                  # Visual Studio 解决方案
    +-- Assets/
    |   +-- Scenes/
    |   |   +-- scene.scene               # 主场景
    |   +-- Scripts/
    |   |   +-- DesktopPet.cs              # 主控 + 物理 + 状态机
    |   |   +-- Live2DRenderer.cs          # Live2D 渲染 + 动画 + 物理驱动
    |   |   +-- ContextMenu.cs             # 右键菜单 UI
    |   |   +-- DragHandler.cs             # 拖拽/抛掷/点击交互
    |   |   +-- WindowOverlay.cs           # Win32 透明窗口
    |   |   +-- TimeWeatherController.cs   # 昼夜/天气响应
    |   |   +-- ChatBubble.cs              # 头顶气泡
    |   |   +-- ChatManager.cs             # AI 对话管理
    |   |   +-- AutoChat.cs                # 自动闲聊
    |   |   +-- BottomInputBar.cs          # 底部输入栏
    |   |   +-- HybridRenderer.cs          # 混合渲染管理器
    |   |   +-- IPetRenderer.cs            # 渲染器接口
    |   |   +-- Model3DRenderer.cs         # 3D 渲染器 (预留)
    |   |   +-- DebugWindow.cs             # 调试调参面板
    |   |   +-- ToolCallInvoker.cs         # AI 工具调用
    |   +-- Live2D/
    |   |   +-- Models/Fuxuan/             # 符玄 Live2D 模型
    |   +-- Plugins/                       # Cubism SDK
    +-- Packages/
    |   +-- manifest.json                  # Unity 包管理
    +-- ProjectSettings/                   # Unity 项目设置
\end{lstlisting}

% ═══════════════════════════════════════════
\section{开发指南}
% ═══════════════════════════════════════════

\subsection{调试快捷键}

\begin{itemize}
    \item \texttt{1}：强制切换到 Live2D 渲染（HybridRenderer 调试用）
    \item \texttt{2}：强制切换到 3D 渲染（HybridRenderer 调试用）
    \item 控制台日志：各组件关键事件均有 Debug.Log 输出
\end{itemize}

\subsection{调参指南}

\begin{itemize}
    \item \textbf{模型大小/位置}：修改 Live2DRenderer.cs 顶部的
        \texttt{LIVE2D\_SCALE}（默认 200）和 \texttt{LIVE2D\_OFFSET\_Y}（默认 100）
    \item \textbf{动作时长}：各动作对应的 \texttt{*\_DURATION} 常量
    \item \textbf{权重调整}：通过右键菜单 UI 实时调整
    \item \textbf{过渡平滑度}：\texttt{WALK\_FADE\_IN\_DURATION}（0.3s）、
        \texttt{IDLE\_BLEND\_DURATION}（0.4s）
    \item \textbf{物理}：DesktopPet.cs 顶部的重力/速度常量
\end{itemize}

\subsection{构建注意}

\begin{itemize}
    \item Camera Background 必须为 \#00FF00（色键透明色）
    \item 构建版自动应用透明窗口设置（编辑器版不执行）
    \item Live2D 模型需通过 Cubism SDK 导入生成 Prefab
    \item 单例互斥锁防止 Build \& Run 产生多个实例
\end{itemize}

% ═══════════════════════════════════════════
\section{版本历史}
% ═══════════════════════════════════════════

\begin{table}[H]
\centering
\caption{版本历史}
\label{tab:history}
\begin{tabular}{cl>{\raggedright\arraybackslash}p{0.55\textwidth}}
\toprule
\textbf{日期} & \textbf{版本} & \textbf{变更} \\
\midrule
2025 & v0.1 & 透明窗口 + 基础物理 + Live2D 加载 \\
2025 & v0.2 & 地面状态机 + 走路动画 \\
2025 & v0.3 & 10 种空闲动作 + 强制动作系统 \\
2025 & v0.4 & 右键菜单 UI + 权重编辑 \\
2025 & v0.5 & 动作锁定机制 + 走路淡入过渡 \\
2026-06-13 & v0.6 & 本文档 + 项目文档整理 \\
2026-06-17 & v0.7 & 修复行走手臂被物理覆盖问题：添加 \texttt{[DefaultExecutionOrder(801)]} \\
           &      & 确保 LateUpdate 跑在 CubismPhysicsController 之后；\\
           &      & 体态参数提前到 Update() 设置供物理驱动衣服；\\
           &      & 分离大/小范围手臂参数常量避免 clamp 满幅；\\
           &      & 走路时 ForceUpdateNow 强制网格重算 \\
2026-06-17 & v0.8 & 头发/裙子/法盘直接物理驱动（20+ 输出参数绕过物理延迟）\\
           &      & 分区点击反馈（头/身/腿 3 区 + 参数保存防覆盖）\\
           &      & 拖拽挣扎动画（双臂划水 + 双腿交替 + 扭动 + 表情）\\
           &      & 时间/天气响应系统（wttr.in API + 昼夜犯困 + 天气表情）\\
           &      & 新增第 11 个空闲动作「困惑」\\
           &      & 待机气泡支持时间/天气特化内容\\
           &      & 更新项目文档 README + LaTeX\\
\bottomrule
\end{tabular}
\end{table}

\end{document}

检测鼠标点击模型时的反馈逻辑。

\begin{lstlisting}[style=csharp,caption={TouchController.cs 骨架},label={lst:touch}]
using UnityEngine;

public class TouchController : MonoBehaviour
{
    [Header("点击反馈")]
    public float reactionDuration = 1.0f; // 反应持续时间

    private DesktopPet pet;
    private float reactionTimer = 0f;
    private bool isReacting = false;

    void Start()
    {
        pet = GetComponent<DesktopPet>();
    }

    void OnMouseDown()
    {
        if (!isReacting)
        {
            // 触发点击反应
            isReacting = true;
            reactionTimer = reactionDuration;
            pet.SwitchState(DesktopPet.PetState.Idle);

            // 随机播放一个反应动画
            Animator anim = GetComponent<Animator>();
            if (anim != null)
                anim.SetTrigger("Touch");

            // 可选：播放音效
        }
    }

    void Update()
    {
        if (isReacting)
        {
            reactionTimer -= Time.deltaTime;
            if (reactionTimer <= 0f)
            {
                isReacting = false;
                // 恢复正常状态
            }
        }
    }
}
\end{lstlisting}

% ═══════════════════════════════════════════
\section{行走动画设计原理}
% ═══════════════════════════════════════════

\noindent
\textbf{版本}：1.1 \\
\textbf{更新内容}：补充 2D 侧面行走循环（Walk Cycle）动画理论，指导 Live2D 参数化走路动画的实现。

\subsection{核心哲学}

行走的本质是\textbf{「摔倒—接住自己」}的循环：
重力不断将身体拉向地面，而双腿交替前伸，在身体即将倾倒的瞬间接住它。
这一个循环中提炼出 4 个关键姿势（Key Poses），它们定义了行走的节奏和力学感。

\subsection{四大关键姿势}

\begin{table}[H]
\centering
\caption{行走循环四大关键姿势}
\label{tab:walkcycle}
\begin{tabular}{lp{3.5cm}p{6cm}c}
\toprule
\textbf{姿势} & \textbf{别名} & \textbf{描述} & \textbf{身体高度} \\
\midrule
Contact & 跨步位、Heel \& Toe & 一脚跟触地、一脚尖离地，\newline 身体处于「失重」瞬间 & 低 \\
Down / Recoil & 下落位、重量转移 & 重心完全压在前脚，\newline 膝盖弯曲吸收冲击 & \textbf{最低} \\
Passing & 过地位 & 承重腿在身体正下方，\newline 另一条腿经过其旁 & \textbf{最高} \\
Up / High & 抬高位、升空位 & 后腿推离地面，\newline 前腿伸展准备落地 & 高 \\
\bottomrule
\end{tabular}
\end{table}

\subsection{垂直弹跳相位}

垂直弹跳是行走动画中\textbf{最重要}的视觉信号——没有它，行走就会变成「地面滑行」。

\begin{itemize}
    \item \textbf{最高点}：Passing（过地位）——小腿肌肉将身体推起
    \item \textbf{最低点}：Down（下落位）——膝盖弯曲吸收冲击
    \item \textbf{相位序列}：Contact(低)→Down(\textbf{最低})→Passing(\textbf{最高})→Up(高)→Contact(对侧)
    \item \textbf{弹跳幅度}：自然行走约身高的 2--4\%，卡通可放大到 10--20 cm
\end{itemize}

数学表达：设步态相位 $\phi \in [0, 2\pi)$，则

\begin{equation}
    \text{BodyHeight}(\phi) = A \cdot \cos(\phi) + \text{BaseHeight}
\end{equation}

其中 $A$ 为弹跳幅度，$\phi=0$ 时（Passing 位）身体最高，$\phi=\pi$ 时（Down 位）身体最低。

\subsection{手臂与腿的交叉对位}

\textbf{核心法则}：一条腿向前时，\textbf{对侧}手臂向前（同侧手臂向后）。

\begin{equation}
    \text{Arm}_{\text{left}} = -\text{Leg}_{\text{right}} = +\text{Leg}_{\text{left}}
\end{equation}

\begin{table}[H]
\centering
\caption{手臂-腿交替关系}
\label{tab:armleg}
\begin{tabular}{cccc}
\toprule
\textbf{阶段} & \textbf{左腿} & \textbf{左臂} & \textbf{右臂} \\
\midrule
左腿在前 (Contact) & 前 & 后 & 前 \\
左腿过地 (Passing) & 中 & 中 & 中 \\
右腿在前 (Contact) & 后 & 前 & 后 \\
\bottomrule
\end{tabular}
\end{table}

这一模式来自骨盆与肩膀的\textbf{对向旋转}，是人体自然步态的生物学基础。

\subsection{身体前倾与速度}

\begin{table}[H]
\centering
\caption{行走类型与前倾角度}
\label{tab:lean}
\begin{tabular}{lc}
\toprule
\textbf{行走类型} & \textbf{前倾角度} \\
\midrule
休闲散步 & 3--7° \\
正常步行 & 5--10° \\
快走 & 10--20° \\
奔跑 & 20--30°+ \\
\bottomrule
\end{tabular}
\end{table}

前倾角度与速度呈正比：速度越快，需要更大的前倾来产生水平方向的动量。

\subsection{头部运动}

\begin{itemize}
    \item \textbf{随身体弹跳}：头上下的幅度约为身体弹跳的 50--70\%
    \item \textbf{平滑补偿}：大脑为保持视野水平，会使头部轨迹比骨盆更平滑
    \item \textbf{二次运动}：头部可做微小延迟跟随（跟随主运动后到达），增加生动感
    \item \textbf{行进方向}：朝行走方向微倾，自然行走时约低头 3--5°
\end{itemize}

\subsection{Live2D 参数映射}

结合符玄 Live2D 模型的实际可用参数，将行走循环映射如下：

\begin{table}[H]
\centering
\caption{行走动画 $\to$ Live2D 参数映射}
\label{tab:mapping}
\begin{tabular}{llc}
\toprule
\textbf{动画要素} & \textbf{Live2D 参数} & \textbf{典型值} \\
\midrule
身体转体（侧面视角） & ParamBodyAngleY & 18°（恒定） \\
身体前倾 & ParamBodyAngleX & 5°（恒定） \\
脸转向侧面 & ParamAngleY & 18°（恒定） \\
低头看路 & ParamAngleX & 8°（恒定） \\
垂直颠簸 & 模型 Y 轴（像素偏移） & 4 px \\
呼吸加深 & ParamBreath & 3.0（恒定偏大） \\
\addlinespace
\textbf{左腿参数} & & \\
\quad 抬腿 & Param165 (抬腿) & $\pm$4° \\
\quad 前后摆动 & Param126 (body X 腿 动画 位移) & $\pm$6° \\
\quad 弯曲透视 & Param127 (body X 腿 动画 透视) & 0--6° \\
\addlinespace
\textbf{右腿参数} & & \\
\quad 抬腿 & Param164 (腿) & $\pm$4°（反向） \\
\quad 前后摆动 & Param129 (body X 腿 动画 R 位移) & $\pm$6°（反向） \\
\quad 弯曲透视 & Param131 (body X 腿 动画R 透视) & 0--6°（反向） \\
\addlinespace
\textbf{右臂（大范围）} & Param94 (右手上臂旋转) & $\pm$2（范围[-30,60]） \\
\textbf{右臂R1（前臂）} & Param31 (右手臂R1) & $\pm$0.28（范围[-1,1]） \\
\textbf{右臂R2} & Param32 (右手臂R2) & $\pm$0.16（范围[-1,1]） \\
\textbf{右臂上臂} & Param33 (右手臂R1) & $\pm$0.16（范围[-1,1]） \\
\textbf{左臂L1} & Param34 (手臂L1) & $\pm$0.28（范围[-1,1]，反向） \\
\textbf{左臂L2} & Param36 (手臂L2) & $\pm$0.16（范围[-1,1]，反向） \\
\textbf{左臂L3} & Param37 (手臂L3) & $\pm$0.16（范围[-1,1]，反向） \\
耸肩 & Param153 & 0--1.5° \\
\bottomrule
\end{tabular}
\end{table}

\subsection{常见错误}

\begin{enumerate}
    \item \textbf{身体高度无变化}（滑动步态 / Sliding）—— 最常见的错误，行走看起来像在地面滑行
    \item \textbf{无重量转移} —— 重心始终在中间，缺乏脚踏实地感
    \item \textbf{同手同脚} —— 混淆了交叉对位法则
    \item \textbf{弹跳相位错误} —— 最高点和最低点位置搞反，导致颠簸节奏违和
    \item \textbf{头部刚性跟随} —— 头应该比身体有轻微的延迟和独立运动
    \item \textbf{速度与参数不协调} —— 步频、步幅、前倾角度之间应保持一致性
\end{enumerate}

% ═══════════════════════════════════════════
\section{实施路线}
% ═══════════════════════════════════════════

项目实施分为四个阶段，每阶段有明确的里程碑：

\subsection{第一阶段：环境搭建（第 1 天）}

\begin{itemize}
    \item 在 Unity Hub 中创建新项目 \texttt{DesktopPerPro}
    \item 导入 \texttt{TransparentWindowManager.unitypackage}
    \item 配置 Camera 为透明背景
    \item 验证：运行后 Unity 窗口透明叠加在桌面上
    \item 初始化 Git 仓库，创建 \texttt{.gitignore}（Unity 模板）
\end{itemize}

\subsection{第二阶段：模型导入（第 2--3 天）}

\begin{itemize}
    \item 获取符玄 3D 模型（FBX/Unity AssetBundle）
    \item 导入 Unity，配置贴图和材质
    \item 建立 Animator Controller，设置基本动画状态
    \item 搭建简单场景：将模型放置在透明背景中
\end{itemize}

\subsection{第三阶段：交互实现（第 4--6 天）}

\begin{itemize}
    \item 实现 \texttt{DesktopPet.cs}：状态机 + 物理
    \item 实现 \texttt{DragHandler.cs}：拖拽 + 抛掷
    \item 实现 \texttt{WindowOverlay.cs}：透明 + 置顶
    \item 实现 \texttt{TouchController.cs}：点击反馈
    \item 测试：宠物可在桌面自由拖拽、下落、行走
\end{itemize}

\subsection{第四阶段：集成与优化（第 7--10 天）}

\begin{itemize}
    \item 集成进程间通信（与旧 \texttt{text1.exe} 对话）
    \item 添加窗口置顶轮询（防止被遮挡）
    \item 性能优化（降低无关渲染开销）
    \item 打包为独立 Windows 可执行文件
\end{itemize}

% ═══════════════════════════════════════════
\section{模型获取方案}
% ═══════════════════════════════════════════

由于米哈游官方未直接公开角色模型文件，获取符玄 3D 模型有以下途径：

\begin{enumerate}
    \item \textbf{SR-Model-Importer}：社区开发的 Unity 工具，可从游戏资源包中提取角色模型
    \item \textbf{hoyo-models 社区}：部分站点提供已导出的 FBX 格式
    \item \textbf{VRM 社区模型}：有爱好者制作的 VRM 版本，可转换为 FBX
    \item \textbf{替代方案}：先用 Unity Asset Store 的免费角色模型调试交互逻辑，
          待模型到位后再替换
\end{enumerate}

\begin{quote}
\small
\textbf{注意}：以上途径获取的模型仅限个人学习和研究使用，请勿用于商业用途或公开发布。
\end{quote}

% ═══════════════════════════════════════════
\section{Live2D 过渡预留}
% ═══════════════════════════════════════════

为将来迁移至 Live2D 技术，在本项目中预留以下扩展点：

\begin{itemize}
    \item \textbf{抽象渲染接口}：将宠物渲染抽象为接口 \texttt{IPetRenderer}，
          当前实现为 3D 模型渲染，后续可替换为 Live2D Cubism SDK 渲染
    \item \textbf{参数化系统}：当前使用 Unity Animator 的参数（Float/Trigger），
          与 Live2D 的参数系统（ParamAngleX/ParamEyeLOpen 等）概念一致，
          便于映射
    \item \textbf{交互逻辑复用}：拖拽、点击、物理等交互逻辑与渲染方式解耦，
          更换渲染引擎时无需重写交互代码
\end{itemize}

\begin{lstlisting}[style=csharp,caption={渲染接口抽象（预留）},label={lst:interface}]
// 预留的渲染接口，方便后续替换为 Live2D
public interface IPetRenderer
{
    void SetParameter(string name, float value);
    void PlayAnimation(string animationName);
    void SetExpression(string expressionName);
    void SetMouthOpen(float value);    // 说话口型
    void SetEyeOpen(float left, float right); // 眨眼
}

// Unity 3D 实现
public class UnityModelRenderer : MonoBehaviour, IPetRenderer
{
    private Animator animator;
    private SkinnedMeshRenderer meshRenderer;

    public void SetParameter(string name, float value)
    {
        if (animator != null)
            animator.SetFloat(name, value);
    }

    public void SetEyeOpen(float left, float right)
    {
        // 通过 Blend Shape 控制眼睛
    }
    // ... 其他接口实现
}

// 未来: Live2D 实现
// public class Live2DRenderer : MonoBehaviour, IPetRenderer { ... }
\end{lstlisting}

% ═══════════════════════════════════════════
\section{目录结构规划}
% ═══════════════════════════════════════════

\begin{verbatim}
Desktop_per_pro/              # Git 仓库根目录
  |-- README.md               # 项目说明
  |-- report.pdf              # 本文档（可删除）
  |-- report.tex              # 本文档源码（可删除）
  |-- Assets/
  |     |-- Scenes/
  |     |     +-- Main.unity  # 主场景
  |     |-- Scripts/
  |     |     |-- DesktopPet.cs
  |     |     |-- DragHandler.cs
  |     |     |-- WindowOverlay.cs
  |     |     |-- TouchController.cs
  |     |     |-- IPetRenderer.cs     # 渲染接口
  |     |     +-- UnityModelRenderer.cs
  |     |-- Models/
  |     |     +-- FuXuan/             # 符玄模型
  |     |-- Materials/
  |     |-- Animations/
  |     |     +-- AnimatorController.controller
  |     +-- Plugins/
  |           +-- TransparentWindowManager/
  |-- Packages/
  |-- ProjectSettings/
  +-- .gitignore
\end{verbatim}

% ═══════════════════════════════════════════
\begin{thebibliography}{99}
% ═══════════════════════════════════════════

\bibitem{TransparentWindow}
XJINE. \textit{Unity\_TransparentWindowManager --- Make Unity's window transparent and overlay on desktop}.
\url{https://github.com/XJINE/Unity_TransparentWindowManager}, 2018. MIT License.

\bibitem{HoyoModels}
hoyo-models Community. \textit{Honkai: Star Rail 3D Character Models}.
\url{https://hoyo-models.vercel.app/hsr}, 2025.

\bibitem{SRImporter}
lethern. \textit{SR-Model-Importer --- Honkai: Star Rail Model Importer for Unity}.
\url{https://github.com/lethern/SR-Model-Importer}, 2025.

\bibitem{PetOriginal}
原桌面宠物项目. \textit{Desktop Pet --- GDI+ + WebView2 桌面宠物}.
路径：\texttt{D:/C/Desktop pet/text1.c}, 2025.

\end{thebibliography}

\end{document}
